% \chapter{绪论}\label{chap:introduction}
\begin{cuzchapter}{绪论}{chap:introduction}

	%---------------------------------------------------------------------------%
	%->> 绪论写作指南
	%---------------------------------------------------------------------------%
	% 绪论是论文的第一章，主要介绍研究的背景、意义、目的、内容和方法等
	% 一个完整的绪论应包含以下几个部分：
	% 1. 研究背景：介绍研究领域的发展历程和现状，指出存在的问题
	% 2. 研究意义：阐述开展本研究的理论意义和实践价值
	% 3. 研究目的：明确说明本研究要解决的具体问题
	% 4. 研究内容：概述论文的主要研究内容和章节安排
	% 5. 研究方法：简要介绍本研究采用的主要研究方法

	\section{研究背景}\label{sec:background}

	如今，幻灯片是表达演讲者想法和便利沟通的重要工具，传统的幻灯片工具如Microsoft PowerPoint和Apple Keynote采用“所见即所得”的编辑方式，
	虽然这种方式在许多地方符合直觉且容易上手，但在团队协作、版本控制、代码展示等方面仍然存在许多不便。
	随着远程协作的兴起，我们需要支持多人实时协作、版本管理和智能辅助等功能的幻灯片制作工具。
	% 如今，幻灯片演示成为教育、商业、科研等领域中不可或缺的沟通工具。传统的幻灯片工具如Microsoft PowerPoint和Apple Keynote采用所见即所得的编辑方式，虽然在许多地方符合直觉且容易上手，但在团队协作、版本控制、代码展示等方面仍然存在许多不便。随着远程协作的兴起，我们需要支持多人实时协作、版本管理和智能辅助等功能的幻灯片制作工具。

	近年来，基于Web的轻量化幻灯片工具\cite{XXDL201613117}逐渐出现在人们的视野，
	如Reveal.js、Marp、Slidev等的工具支持用Markdown语法编写幻灯片内容，还可以结合现代Web技术栈实现许多有趣且丰富的交互效果。
	其中，Slidev（slide + dev）作为一个专门为开发者设计的现代幻灯片制作工具，不仅支持Vue语法，而且支持丰富的自定义插件，让越来越多开发者开始关注并使用它了。


	综合以上背景，作者设计与开发一个基于Slidev的智能幻灯片协作平台，支持实时协作技术、权限管理系统、AI辅助功能和简化版本控制。

	\section{研究意义}\label{sec:significance}

	开发基于Slidev的智能幻灯片协作平台有着重要的理论和实践意义：

	\subsection{理论意义}

	\textbf{探寻合适的人机交互方式}：摸索大语言模型在幻灯片内容生成与优化中的应用，探寻人机协作的较为合适的交互方式。
	% \textbf{探索AI辅助创作模式}：研究大语言模型在幻灯片内容生成与优化中的应用，探索人机协作的智能创作模式，为AI辅助内容创作提供理论参考。
	\textbf{研究权限管理模型}：基于RBAC\cite{zhang2013rbac}（基于角色的访问控制）模型，设计适用于幻灯片多人协作场景的权限管理系统。

	\subsection{实践意义}

	\textbf{提高团队协作效率}：支持多个用户实时协同编辑同一幻灯片，而不需要来来回回地传输文件，以此提高团队工作效率和减少沟通成本。
	\textbf{降低技术使用门槛}：通过基于快照的版本管理来避免Git\cite{SXDS202505026}等专业工具，降低了非技术用户的使用门槛。
	 \textbf{助力内容创作}：利用加入Slidev相关知识库的AI来帮助用户快速生成幻灯片大纲、优化内容。
	%  \textbf{增强内容创作能力}：集成AI辅助功能，帮助用户快速生成幻灯片内容、设计布局和优化表达，提升内容创作质量和效率。
	  \textbf{方便团队的知识共享与管理}：通过幻灯片空间和文档树形结构，能够很方便地帮助团队组织和管理技术文档。

	\section{研究目的}\label{sec:purpose}

	作者想要设计并实现一个功能完善、方便非技术人员也能快速上手使用的智能幻灯片协作平台，
	大致目标如下
	：开发基于React + NestJS的Web平台，实现Slidev代码编辑器的在线协作；
	支持多用户同时编辑和显示其他协作者正在编辑的光标区域；
	设计基于RBAC的四级权限管理系统（Owner/Editor/Commenter/Viewer），实现合适的文档权限访问控制；
	接入大语言模型API，辅助幻灯片内容生成、语法适配和提供智能建议；
	实现简单的、基于快照的版本控制系统，提供版本历史查看、差异对比和回滚功能；
	在幻灯片空间中，支持文档树形结构和标签分类；
	支持对幻灯片的评论，方便团队协作和交流。

	\section{研究内容}\label{sec:content}

	本研究的主要内容包括以下几个方面：

	\subsection{平台架构设计}

	系统采用前端分离架构，
	包括前端展示层、后端业务层和数据存储层。
	前端使用的是React 19、TypeScript和Vite，将CodeMirror 6编辑器\cite{huang2019codemirror}作为Slidev的代码编辑器，并将Yjs作为实时协作框架。
	后端依靠NestJS框架来构建，提供RESTful API\cite{wang2013restful,zhang2018restful}和WebSocket协作服务。
	数据存储层选择MySQL数据库，并通过TypeORM进行数据持久化。

	\subsection{实时协作功能实现}

	用Yjs库实现操作转换和冲突解决。以Hocuspocus WebSocket 服务器来处理多客户端连接，实现文档状态同步和实时光标显示。
	% 研究CRDT算法在文档协作中的应用，基于Yjs库实现操作转换和冲突解决。用Hocuspocus WebSocket 服务器来处理多客户端连接，实现文档状态同步和实时光标显示。
	\subsection{权限管理系统开发}

	基于JWT认证和RBAC模型\cite{lu2025jwtrbac,li2011reengineering}，设计四级用户角色权限体系。实现文档可见性控制（公开/私有）、协作成员管理和操作权限校验，确保文档安全性和访问控制。

	\subsection{AI辅助功能集成}

	集成通义千问等大语言模型API，实现幻灯片内容生成、语法适配、布局建议等功能。通过加入Slidev文档作为AI的知识库，让AI生成的内容更好地适配Slidev语法。

	\subsection{版本控制系统设计}
	
	实现简单的、基于快照的版本控制，支持自动保存和手动保存、手动创建版本，并提供版本历史查看、内容差异对比和一键回滚操作，让非技术用户更快、更容易上手。
	% 设计简化的版本控制流程，实现自动保存和手动版本创建功能。提供版本历史查看、内容差异对比和一键回滚操作，降低非技术用户的使用难度。

	\section{研究方法}\label{sec:methodology}

	本研究采用以下研究方法：
	\textbf{文献研究法}：查阅国内外相关文献，了解幻灯片制作工具、实时协作技术、权限管理和AI辅助创作等领域的研究现状和发展趋势。 
	
	\textbf{系统设计法}：进行需求分析、系统设计、模块划分和接口定义，确保系统架构的合理性和可扩展性。 
	
	\textbf{原型开发法}：采用敏捷开发，迭代实现系统核心功能，通过快速原型验证技术方案是否合适于本系统。 

	\textbf{测试验证法}：通过黑盒测试验证系统核心功能的正确性，确保实时协作、权限管理、AI辅助和版本控制等关键模块符合设计要求。 
	
	\textbf{案例分析法}：选取一些关键或常见的使用场景进行案例研究，分析系统在实际应用中的效果并探寻改进方案。

	\section{论文结构}\label{sec:structure}

	本论文共分为五章，各章内容概述如下：

	\textbf{第一章 绪论}：介绍研究背景、研究意义、研究目的、研究内容和研究方法，阐述开发智能幻灯片协作平台的现实意义和作用。

	\textbf{第二章 文献综述}：回顾国内外相关研究成果，包括幻灯片制作工具、实时协作技术、权限管理系统和AI辅助创作等方面的研究，为平台开发提供理论基础。

	\textbf{第三章 系统设计与实现}：详细介绍平台的设计思路、架构和实现方法，包括技术选型、模块设计、核心算法和关键技术的实现细节。

	\textbf{第四章 使用指南}：介绍平台的使用方法和操作流程，包括平台概述、核心功能使用指南、高级功能详解和最佳实践与技巧，帮助用户快速上手。

	\textbf{第五章 总结与展望}：总结平台开发的主要成果和特点，分析存在的不足，并对未来的改进方向和发展前景进行展望。

	% 此外，论文还包括参考文献、致谢和附录等部分，提供补充信息和资源。
	此外，论文还包括参考文献、致谢等部分，提供补充信息和资源。

\end{cuzchapter}